\pdfvariable suppressoptionalinfo 611
\documentclass[10pt]{article}
\usepackage[a4paper,margin=0.52in]{geometry}
\usepackage{fontspec}
\setmainfont{Latin Modern Roman}
\setsansfont{Latin Modern Sans}
\setmonofont{Latin Modern Mono}
\usepackage[english]{babel}
\babelprovide[import]{chinese-simplified}
\babelfont[chinese-simplified]{rm}[Renderer=Harfbuzz,Language=Default]{Droid Sans Fallback}
\usepackage{amsmath,amssymb}
\linespread{0.92}
\emergencystretch=2em
\title{\vspace{-1.8em}Exact Walsh Traces and Krawtchouk Compression\\for the Ehrenfest Hypercube Walk}
\author{Route-A structural certificate C171}
\date{}
\begin{document}
\maketitle
\begin{abstract}
For every integer $d\geq1$, we diagonalize the Ehrenfest coordinate-flip
operator on the $d$-dimensional sign cube.  Every Walsh character indexed by
$S$ has eigenvalue $1-2|S|/d$, with binomial multiplicity.  This gives exact
power traces, a factored determinant, and every vertex return probability;
odd returns vanish by bipartiteness.  Hamming weight is an exact Markov
lumping to a $(d+1)$-state reversible birth--death chain.  Its symmetric
Jacobi similarity has off-diagonal entries
$\sqrt{(k+1)(d-k)}/d$, and Krawtchouk polynomials provide its simple spectrum,
retaining every distinct eigenvalue of the original $2^d$-state operator.
The natural $P_d$ is self-adjoint, but exponentiation changes the frozen
Markov clock and weights.  The family has no intrinsic prime correspondence;
the result is an all-parameter source theorem, not a target-spectrum claim.
\end{abstract}
\noindent\textbf{Keywords:} Ehrenfest chain; hypercube dynamics; Walsh
characters; Krawtchouk polynomials; reversible Markov chains; exact traces.

\begin{otherlanguage}{chinese-simplified}
\begin{center}{\large 中文摘要}\end{center}
\begin{sloppypar}
对任意整数 $d\geq1$，本文精确对角化 $d$ 维符号超立方体上的
\foreignlanguage{english}{Ehrenfest} 坐标翻转算子。由集合 $S$ 标记的
\foreignlanguage{english}{Walsh} 特征函数具有特征值 $1-2|S|/d$，其重数为二项式系数。
由此得到精确幂迹、因子化行列式与任意顶点的返回概率；二部结构使所有奇数时刻返回概率为零。
汉明重量给出到 $d+1$ 状态可逆生灭链的精确马尔可夫压缩。其对称
\foreignlanguage{english}{Jacobi} 相似矩阵的非对角元为
$\sqrt{(k+1)(d-k)}/d$，而 \foreignlanguage{english}{Krawtchouk} 多项式给出单谱，
从而保留原 $2^d$ 状态算子的每个不同特征值。自然算子 $P_d$ 虽然自伴，但取酉指数会改变冻结的
马尔可夫时钟与路径权重；本系统也不内生素数对应，因此本文不作目标谱结论。
\end{sloppypar}
\noindent 关键词：\foreignlanguage{english}{Ehrenfest} 链；超立方体动力学；
\foreignlanguage{english}{Walsh} 特征；\foreignlanguage{english}{Krawtchouk} 多项式；可逆马尔可夫链；精确迹。
\end{otherlanguage}

\section{Frozen dynamics}
Let $X_d=\{-1,+1\}^d$, let $x^{(i)}$ flip coordinate $i$, and set
\[
 (P_df)(x)=\frac1d\sum_{i=1}^d f(x^{(i)})
 \quad\text{on }L^2(X_d,2^{-d}\#).                              \tag{1}
\]
One coordinate flip is one Markov tick.  For $S\subseteq\{1,\ldots,d\}$,
write $\chi_S(x)=\prod_{i\in S}x_i$.

\paragraph{Theorem 1 (complete Walsh certificate).}
For every $d\geq1$, the Walsh characters form an orthonormal eigenbasis and
\[
 P_d\chi_S=\lambda_{|S|}\chi_S,\qquad
 \lambda_j=1-\frac{2j}{d},\qquad
 \operatorname{mult}(\lambda_j)=\binom dj.                     \tag{2}
\]
Consequently, for every $n\geq0$,
\begin{align}
 \operatorname{Tr}P_d^n&=\sum_{j=0}^d\binom dj\lambda_j^n,       \tag{3}\\
 \det(I-zP_d)&=\prod_{j=0}^d(1-z\lambda_j)^{\binom dj}.          \tag{4}
\end{align}

\paragraph{Proof.}
A flip contributes $-\chi_S$ for the $|S|$ coordinates in $S$ and
$+\chi_S$ for the other $d-|S|$.  Averaging proves (2).  Walsh orthogonality
and the count $\binom dj$ give a complete basis; (3)--(4) follow from the
finite spectral theorem. \hfill$\square$

For $|z|<1$, the scalar power series for $-\log(1-z\lambda_j)$ gives the
exact trace--log relation
\[
 -\log\det(I-zP_d)=\sum_{n\geq1}\operatorname{Tr}(P_d^n)\frac{z^n}{n}.
                                                                    \tag{5}
\]
For every $d>1$ this is a weighted closed-Markov-walk identity: $P_d$ is a
genuine average rather than a deterministic map, so these traces are not
fixed-point counts and (5) is not an Artin--Mazur zeta.  At $d=1$, $P_1$ is
the deterministic two-cycle permutation; that isolated boundary does not
supply a uniform all-$d$ primitive-orbit interpretation.

\section{Return law and exact compression}
Cube translations commute with $P_d$ and act transitively, so all diagonal
entries of $P_d^n$ agree.  Their sum is the trace, hence
\[
 P_d^n(x,x)=2^{-d}\operatorname{Tr}P_d^n.                       \tag{6}
\]
Every flip reverses Hamming parity, proving that (6) vanishes for odd $n$.

Let $k$ count negative coordinates.  The exact lumped kernel is
\[
 Q(k,k+1)=\frac{d-k}{d},\qquad Q(k,k-1)=\frac{k}{d}.             \tag{7}
\]
For $\pi_k=2^{-d}\binom dk$, the binomial identity
$\binom dk(d-k)=\binom d{k+1}(k+1)$ proves detailed balance.  Therefore
$D_\pi^{1/2}QD_\pi^{-1/2}$ is symmetric with upper entry
$\sqrt{(k+1)(d-k)}/d$.

\paragraph{Theorem 2 (Krawtchouk compression).}
Define
\[
 K_j(k)=\sum_r(-1)^r\binom kr\binom{d-k}{j-r}.
\]
Then $QK_j=(1-2j/d)K_j$.  Thus $Q$ has the $d+1$ distinct eigenvalues
in (2), each once, and compresses the $2^d$-state operator without losing a
distinct eigenvalue.

\paragraph{Proof.}
The generating function is
$\sum_jK_j(k)t^j=(1-t)^k(1+t)^{d-k}$.  Substitution of $k\pm1$ using (7)
equals the original series minus $(2t/d)$ times its derivative.  Comparing
$t^j$ coefficients gives the eigenvalue equation.  The eigenvalues are
distinct, so the $K_j$ form a basis. \hfill$\square$

\section{Arithmetic and operator controls}
Four frozen controls expose the scope.  Random arithmetic labels do not enter
$P_d$; composite-only dimension labels change no theorem; neighboring
dimensions obey the same binomial law; and $\eta I+(1-\eta)P_d$ has the
affinely shifted spectrum for any source-chosen $\eta$.  For $d>1$ there is no
deterministic primitive-orbit owner.  The $d=1$ two-cycle is an isolated
boundary with no all-family repetition law, $\log p$ clock, or arithmetic
weight.  Thus A0 fails and A1 fails for the full family.

Each coordinate flip is a self-adjoint permutation, so their average $P_d$ is
a natural self-adjoint contraction.  However, $e^{-itP_d}$ is not the frozen
one-flip Markov dynamics: it changes both clock and path weights.  This is only
a formal A4 hint, not a same-clock quantization.  The Route-A tuple is
\[
 (\mathrm{A0\_FAIL},\mathrm{A1\_FAIL},\mathrm{A2\_FAIL},
   \mathrm{A3\_FAIL},\mathrm{A4\_FORMAL\_HINT}).                 \tag{8}
\]
Because A0 is mandatory, the overall Route-A status is
\texttt{ROUTE\_A\_REJECTED}; the exact source theorem is retained.

\section{Boundary}
For $d=1$, the theorem reduces to the deterministic two-cycle with spectrum
$\{1,-1\}$.  At even $d$, zero has multiplicity $\binom d{d/2}$; its factor
in (4) is one and contributes no polynomial degree.  At $n=0$, (3) and (6)
give $2^d$ and one, respectively.  The unit eigenvalue makes $|z|<1$ the
natural domain stated for (5).

Exact ledgers through $d=18,n=24$, independent reconstruction, brute
closed-walk enumeration for $d\leq7,n\leq8$, symbolic Krawtchouk checks, byte
replay and repaired-hash mutations test the implementation; they do not prove
the theorem.  No target divisor, functional equation, counting law,
arithmetic local datum, Euler factor, root number, automorphy,
Hilbert--Polya operator or Route-B authorization is asserted.  The scope is
\texttt{NO\_BAD\_EULER\_OR\_ROOT\_NUMBER}.

\paragraph{Declarations.}
All data and code needed to reproduce the exact artifact are included.  No
human or animal subjects, private data or external dataset are involved.  The
anonymous contribution record covers conceptualization, formal analysis,
software, validation and writing.  No funding is declared.  No competing
interest is declared.  AI assistance was used for drafting and code
generation; every claim-bearing identity was deterministically reconstructed,
and no external reviewer or acceptance score is represented.
\end{document}
